A ocorrência de enchentes é frequentemente associada a chuvas intensas e ao transbordamento de rios. No entanto, um
fator subjacente e cada vez mais relevante para a intensificação desses eventos é o aquecimento global. O aumento das
temperaturas globais altera o ciclo hidrológico, intensificando a evaporação, modificando os padrões de precipitação e
contribuindo para eventos extremos mais frequentes e intensos. Dessa forma, compreender a relação entre o aquecimento
global e as enchentes é essencial para o desenvolvimento de estratégias eficazes de mitigação e adaptação

\subsection{Aumento de temperatura e emissões globais}
% V. Aumento de temperaturas globais; V. Ação humana no clima; V. Sul Global mais afetado proporcionalmente
% 3. Aumento de heat waves -> Dunn; V. Aumento de precipitação;
% 5. Aumento de Enchentes (precipitação/precipitação acumulada);
% 6. El Nino e La Nina

% Chapter 2.3.1.1.3 (Temperatures during the instrumental period - surface) of the AR6 WGI;
As últimas décadas foram marcadas por quebras constantes de recordes climáticos, especialmente em máximas de
temperatura, evidenciando de forma empírica o impacto das mudanças climáticas em escala global. Esta tendência pode ser
evidenciada através das Anomalias de Temperatura que são definidas como a diferença entre a temperatura observada em um
determinado período e uma temperatura média de referência calculada para um período base. As mesmas são amplamente
utilizadas como indicadores de tendências e mudanças de longo prazo nos padrões climáticos, pois podem capturar
variações consistentes e sistemáticas não aparentes em dados absolutos, permitindo uma análise mais clara das mudanças
globais. Nas últimas décadas, as Anomalias de Temperatura apresentam um padrão de crescimento constante, como pode
evidenciar-se na Figura~\ref{temp_anomaly}.

\begin{figure}[htb]
	\caption{\label{temp_anomaly} Anomalia de Temperatura referente ao período de 1850 a 2024}
	\begin{center}
	    \includegraphics[scale=0.455]{img/Temperature Anomaly.png}
	\end{center}
	\legend{Fonte: \textit{National Oceanic and Atmospheric Administration} (NOAA) \cite{ncei_noaa}}
\end{figure}

De acordo com o sexto - e mais recente - relatório do Painel Intergovernamental sobre Mudanças Climáticas (IPCC),
estabelecido pelas Nações Unidas e amplamente reconhecido como a principal referência na análise do conhecimento
científico relacionado às mudanças climáticas, a superfície terrestre tem sofrido um aumento expressivo de temperatura
em comparação com períodos históricos. Considerando o intervalo de 1850 a 1900 (utilizado como base pré-industrial por
ser o primeiro período com boas observações de temperaturas registradas), a temperatura média da superfície
terrestre\footnote{Dados referentes à medida do GMST - Temperatura Média Global da Superfície - que representa uma média
ponderada entre as temperaturas das massas terrestres e das massas oceânicas.} no período de 2001 a 2020 foi
aproximadamente 1.07ºC mais elevada. Além disso, o aumento da temperatura somente nas últimas quatro décadas (1980 a
2020) é de 0,76ºC, evidenciando uma aceleração no aquecimento global \cite{gulev}.

Esse aumento, aparentemente pequeno, não reflete a real magnitude das modificações no sistema terrestre como um todo. É
importante notar que um aumento de 1ºC representa uma média global; portanto, houveram locais no qual o aquecimento foi
significativamente mais pronunciado, como nos polos. Nessas regiões, mudanças pequenas na temperatura podem levar ao
derretimento de geleiras e, logo, na diminuição do albedo da Terra, isto é, a capacidade de refletir os raios solares de
volta para o espaço. Portanto, nessa situação, a Terra absorverá mais calor, levando à ainda mais derretimento das
geleiras, criando assim um ciclo de retroalimentação. Além disso, é relevante destacar que um aquecimento de 1ºC em toda
a superfície terrestre é extremamente significativo devido à massa do planeta. Para que um aumento dessa magnitude
ocorra em um corpo com a massa da Terra, a quantidade de energia térmica\footnote{De forma bastante simplificada, essa
energia pode ser vista pela fórmula $ Q = mc\Delta T $, onde $m$ é a massa, $c$ é o calor específico, e $\Delta T$ a
variação de temperatura.} liberada é imensa.

% Chapter 3.3 (Human Influence on the Atmosphere and Surface) of the AR6 WGI;
É indubitável que a ação humana foi o maior contribuinte para o Aquecimento Global \cite{ipcc, gillett}. As atividades
industriais que sustentam os padrões de consumo modernos são as principais emissoras de gases responsáveis pelo efeito
estufa, como dióxido de carbono (\ce{CO2}), metano (\ce{CH4}), óxido nitroso (\ce{N2O}), entre outros. Esses gases
possuem a capacidade de absorver a radiação infravermelha emitida pela superfície da Terra, retendo o calor que, de
outra forma, seria dissipado no espaço. 

Embora o efeito estufa seja considerado um fenômeno essencial graças à concentração natural de certos GEEs (sem eles a
temperatura média da superfície terrestre seria significativamente menor), a ação humana tem intensificado esse processo
alarmantemente. Desde o início da Revolução Industrial, a concentração de GEEs na atmosfera aumentou de forma gradual,
principalmente devido à queima de combustíveis fósseis como o carvão, o gás natural e o petróleo. Esses combustíveis são
formados por compostos de carbono de estrutura complexa. Sua alta densidade energética permite que sejam facilmente
queimados para gerar calor, que pode ser convertido em trabalho mecânico em motores, tornando-os a base da matriz
energética global. O principal subproduto desta queima, porém, se dá na forma de gás carbonico (\ce{CO2}) que é liberado
livremente na atmosfera. 

Considerando que a formação de combustíveis fósseis se dá a partir da decomposição de matéria orgânica ao longo de
milhões de anos, a emissão de \ce{CO2} desta forma representa a liberação de carbono na atmosfera que de outra forma
permaneceria preso na superfície. Estima-se que, entre 1850 e 2019, as emissões acumuladas de \ce{CO2} foram de 2400±240
Gt\ce{CO2} \cite{friedlingstein}, resultando em uma concentração de \ce{CO2} na atmosfera de 410 partes por milhão,
comparavél com o início do máximo térmico do Paleoceno-Eoceno \cite{gingerich}, um período associado a uma das mais
recentes extinções em massa.

Esses dados destacam a intensidade das atividades humanas no agravamento do efeito estufa e evidenciam a necessidade
urgente de iniciativas coordenadas em escala global para mitigar os danos climáticos e garantir um futuro mais
sustentável. Iniciativas como o Acordo de Paris são essenciais para garantir a coordenação das ações climáticas ao
estabelecer metas e responsabilidades de cada nação. 

Firmado em 2015 pelas nações participantes do Convenção-Quadro das Nações Unidas sobre a Mudança do Clima (UNFCCC), o
acordo estabelece a necessidade de reduzir as emissões de GEEs em escala global, visando manter o aumento de temperatura
abaixo de 2ºC em relação à base pré-industrial, além de providenciar mecanismos para que nações em desenvolvimento
possam mitigar os efeitos das mudanças climáticas e se adaptar a elas sem comprometer severamente suas atividades
econômicas essenciais.

% Chapter 8.2 of the AR6 WGII and Chapter 12.3.6 of the AR6 WGII;
\subsection{Contribuições e vulnerabilidades climáticas}
\label{vul}
A América Latina é um grande contribuinte na emissão de GEEs, contabilizando aproximadamente 11\% das emissões
históricas globais de \ce{CO2} \cite{friedlingstein}. Contudo, a maior parte dessa emissão vem na forma de uso de terra,
como agricultura, pecuária e desmatamento. Esse padrão é característico de muitas regiões em desenvolvimento, onde a
produção industrial está voltada predominantemente para atender às demandas de exportação de bens e matéria-prima, em
vez de ser direcionada ao consumo interno. Consequentemente, as contribuições dessas regiões para o aquecimento global
vem com o papel de fornecedores globais, beneficiando outras partes do mundo, enquanto enfrentam os desafios associados
a essas emissões.

\begin{figure}[htb]
	\caption{\label{contrib} Emissões Cumulativas por Região (1850-2019)}
	\begin{center}
	    \includegraphics[scale=0.45]{img/contrib_regiao.png}
	\end{center}
	\legend{Fonte: \textit{Global Carbon Budget}, trauduzido pelo autor \cite{friedlingstein}}
\end{figure}

Além disso, essas regiões são significativamente mais vulneráveis aos impactos das mudanças climáticas \cite{chambers},
não apenas por estarem localizadas em áreas tipicamente mais tropicais, onde o aumento das temperaturas é mais intenso,
mas também pela infraestrutura limitada para prevenir e mitigar desastres. Diversos fatores contribuem para a
resiliência de uma região diante das mudanças climáticas, como o desenvolvimento desigual, a fragilidade das
instituições e os altos índices de pobreza. Comunidades mais vulneráveis acabam sendo desproporcionalmente afetadas por
perdas no setor agrícola, o que intensifica a insegurança alimentar e agrava as condições socioeconômicas da região. As
mudanças climáticas também afetam a distribuição dos biomas e sua biodiversidade, criando desafios para as populações
nativas que dependem diretamente desses recursos para sua subsistência e que já vivem em condições de vulnerabilidade.

Na floresta amazônica, por exemplo, o aumento das queimadas, embora frequentemente associado a problemas socioeconômicos
e à negligência governamental, tem sido intensificado pelas mudanças climáticas. O aumento da temperatura e da área
suscetível ao fogo coloca em risco uma das regiões mais biodiversas do planeta e compromete a capacidade da Amazônia de
atuar como um importante sumidouro de carbono, essencial para mitigar as emissões globais. As populações nativas dessas
regiões, apesar de estarem entre as que menos contribuem para as emissões de carbono, são profundamente impactadas pelos
efeitos das queimadas, que aumentam a incidência de doenças respiratórias, e tem a sua mobilidade restringida por
mudanças no ciclo hidrológico que rege o transporte ribeirinho, o principal da região \cite{pinho}. A perda de
biodiversidade afeta diretamente a subsistência desses grupos, que dependem dos recursos naturais para sua alimentação e
sobrevivência. Como consequência da insegurança alimentar e da degradação de seus territórios, muitas comunidades acabam
migrando para centros populacionais próximos, onde enfrentam condições ainda mais vulneráveis, exacerbadas por novos
problemas socioeconômicos e pela dificuldade de adaptação ao ambiente urbano que frequentemente marginaliza suas
necessidades e tradições.

É a intrenseca vulnerabilidade de regiões como esta que torna mais importante as ações de prevenção e mitigação de
desastres naturais causados pelo aumento climático. Pois, mesmo que com a tecnologia atual não se possa reverter o
aumento de temperatura, é possível mitigar os efeitos derivados disto.

% Desenvolvimento sustentável; -> Vulnerabilidade dos que contribuem menos. Tentar ao máximo ficar no tópico de América
% Latina;

% IPCC AR6 WGI Chapter 11.4 (Heavy Precipitation) and Chapter 11.5 (Floods); TODO IMAGEM!
\subsection{Mudanças nos padrões de precipitação}
\label{rain}
Mudanças em média de temperatura são processos graduais que se tornam difícies de perceber em espaços curtos de tempo. A
comparação de temperatura entre um ano e o anterior não mostrará a tendência real das mudanças, especialmente
considerando efeitos naturais de variabilidade climática, como o El Niño e a La Niña. No entanto, essas mudanças
graduais afetam diretamente nos eventos climáticos extremos cuja visibilidade é mais clara. Esse aumento da temperatura,
coloca em desequilibrio delicados mecanismos atmosfericos levando à mudanças em, por exemplo, aumento de temperaturas
máximas, redução de temperaturas mínimas e ondas de calor\footnote{Definido por Dunn et al. (2024) \cite{dunn} como 3 ou
mais dias no qual a temperatura se encontra acima do 90º percentil}.

Um evento climático cujo mudança se destaca é a precipitação. Em essência, o aumento de temperatura leva à uma maior
evaporação da água e também a maior capacidade atmosférica de reter esse vapor de água, o que resulta em chuvas mais
intensas e irregulares. Uma definição mais completa pode ser dada por Seneviratne et al. (2021) no capítulo 11.1 do
sexto relatório, segundo grupo, do IPCC (AR6 WGII): 

\begin{citacao}
As mudanças na temperatura também controlam as variações no vapor d'água por meio do aumento da evaporação e da
capacidade de retenção de água da atmosfera. Em escala global, o conteúdo de vapor d'água integrado na coluna
atmosférica aumenta aproximadamente de acordo com a relação de Clausius–Clapeyron, com um aumento de cerca de 7\% para
cada 1°C de aquecimento médio global da superfície. (...) Projeções de modelos climáticos mostram que o aumento do vapor
d'água resulta em elevações significativas nos extremos de precipitação em todas as regiões, com uma magnitude que varia
entre 4\% e 8\% por 1°C de aquecimento da superfície.

\textit{Seneviratne et al. (2021)\cite{ipcc_ar6_wg2_chap11}, traduzido pelo autor.}

\end{citacao}

(EXPANDIR - Definir  melhor alta precipitação e o aumento dela);

É importante notar que o aumento das temperaturas não está somente associado a um incremento na precipitação total, mas
também à intensificação dos extremos de precipitação. Observa-se não só um aumento no volume máximo de chuva em um único
dia, evidenciando a maior intensidade dos eventos, como também no acumulado de precipitação em períodos de cinco dias,
indicando uma maior duração desses extremos \cite{dunn}.

Apesar da tendência global de aumento nas precipitações extremas, isso não significa que todas as regiões do planeta
experimentarão esse efeito de maneira uniforme. O aumento das chuvas vem acompanhado de mudanças nos padrões de
precipitação, influenciadas pelos mecanismos atmosféricos e pelas características geográficas locais. Assim, o
aquecimento global não apenas eleva o volume de chuvas, mas também altera sua distribuição, afetando onde e com que
frequência ocorrem.

\begin{figure}[htb] 
	\caption{\label{rain_tendency} Tendência de aumento de chuvas em diferentes regiões}
	\begin{center}
	    \includegraphics[scale=0.235]{img/SRL-image-0 (copy).png}
	\end{center}
	\legend{Fonte: \textit{IPCC Sixth Assessment Report - Working Group 1} \cite{ipcc_summary}}
\end{figure}

(EXPANDIR - Mencionar a imagem acima no texto explicando a região SES como alvo);

Essas variações são determinadas por particularidades climáticas regionais, o que dificulta a definição de uma tendência
uniforme. As mudanças nos padrões de precipitação estão diretamente relacionadas à um aumento na quantidade de secas e
de enchentes, pois, se a distribuição de precipitações aumenta ou diminui, os limiares que levam à esses eventos são
mais facilmente alcançados.

No Brasil, por exemplo, esses dois efeitos são observados de forma simultânea. No Nordeste, a região semiárida mais
populosa do mundo, há uma redução na frequência e intensidade das chuvas, o que tem levado a um aumento na ocorrência de
secas mais severas e frequentes, um problema crônico na área. Além disso, projeções indicam que o Nordeste enfrentará
aumentos de temperatura mais intensos do que outras regiões \cite{hoegh-guldberg}, agravando ainda mais os impactos da
seca. Por outro lado, na região Sudeste, observa-se um aumento nas chuvas intensas desde o início do século XX
\cite{dunn}, e, acompanhado das mesmas, uma intensificação em enchentes e deslizamentos, que são a maior causa de morte
relacionada à desastres naturais na região \cite{ufsc}. 
